\documentclass[aspectratio=169,12pt]{beamer}
% Estilo canónico para las presentaciones de Métodos Numéricos 2026-2.
\usepackage[T1]{fontenc}
\usepackage[utf8]{inputenc}
\usepackage[spanish,es-nodecimaldot]{babel}
\usepackage{lmodern}
\usepackage{amsmath,amssymb,mathtools}
\usepackage{booktabs}
\usepackage{graphicx}
\usepackage{xcolor}
\usepackage{hyperref}

\definecolor{upazul}{HTML}{003B70}
\definecolor{upazulclaro}{HTML}{EAF2F8}
\definecolor{updorado}{HTML}{C5962E}
\definecolor{upgris}{HTML}{4A5560}
\definecolor{upverde}{HTML}{167C5A}
\definecolor{uprojo}{HTML}{B33A3A}

\setbeamercolor{normal text}{fg=black,bg=white}
\setbeamercolor{structure}{fg=upazul}
\setbeamercolor{frametitle}{fg=white,bg=upazul}
\setbeamercolor{title}{fg=upazul}
\setbeamercolor{subtitle}{fg=upgris}
\setbeamercolor{block title}{fg=white,bg=upazul}
\setbeamercolor{block body}{fg=black,bg=upazulclaro}
\setbeamercolor{alerted text}{fg=uprojo}
\setbeamercolor{example text}{fg=upverde}

\setbeamerfont{title}{size=\Huge,series=\bfseries}
\setbeamerfont{subtitle}{size=\Large}
\setbeamerfont{frametitle}{size=\LARGE,series=\bfseries}
\setbeamerfont{normal text}{size=\large}
\setbeamerfont{block title}{series=\bfseries}

\setbeamertemplate{navigation symbols}{}
\setbeamertemplate{itemize item}{\color{updorado}$\blacktriangleright$}
\setbeamertemplate{itemize subitem}{\color{upazul}$\bullet$}
\setbeamertemplate{blocks}[rounded][shadow=false]

\setbeamertemplate{frametitle}{%
  \nointerlineskip
  \begin{beamercolorbox}[wd=\paperwidth,ht=0.78cm,dp=0.24cm,leftskip=0.55cm]{frametitle}
    \usebeamerfont{frametitle}\insertframetitle
  \end{beamercolorbox}}

\setbeamertemplate{footline}{%
  \leavevmode
  \hbox{%
    \begin{beamercolorbox}[wd=.82\paperwidth,ht=2.7ex,dp=1.2ex,leftskip=.5cm]{author in head/foot}
      \color{upgris}\scriptsize Métodos Numéricos 2026-2 \hspace{0.8em} | \hspace{0.8em} Universidad del Pacífico
    \end{beamercolorbox}%
    \begin{beamercolorbox}[wd=.18\paperwidth,ht=2.7ex,dp=1.2ex,rightskip=.5cm plus1fil]{date in head/foot}
      \color{upazul}\scriptsize\hfill\insertframenumber/\inserttotalframenumber
    \end{beamercolorbox}}
  \vskip0pt}

\newcommand{\concepto}[1]{\textcolor{upazul}{\textbf{#1}}}
\newcommand{\resultado}[1]{\textcolor{upverde}{\textbf{#1}}}
\newcommand{\advertencia}[1]{\textcolor{uprojo}{\textbf{#1}}}
\newcommand{\codigo}[1]{\colorbox{upazulclaro}{\strut\texttt{#1}}}
\newcommand{\fuente}[1]{\vfill\hfill{\tiny\color{upgris}#1}}


\usepackage{tikz,pgfplots,listings}
\pgfplotsset{compat=1.18}
\lstset{language=Python,basicstyle=\ttfamily\footnotesize,columns=fullflexible,
keepspaces=true,showstringspaces=false,breaklines=true}
\newcommand{\R}{\mathbb R}
\newcommand{\norm}[1]{\left\lVert#1\right\rVert}
\newcommand{\datos}{datos/edo}
\author{Oswaldo José Velásquez Castañón}
\institute{Universidad del Pacífico}
\date{Métodos Numéricos -- 2026-2}
\title{Ecuaciones diferenciales ordinarias}
\subtitle{Problemas de valor inicial y métodos de un paso}
\begin{document}

% 1
\begin{frame}[plain]
\begin{columns}[T,onlytextwidth]
\begin{column}{.73\textwidth}
\vspace{.55cm}{\usebeamerfont{title}\color{upazul}\inserttitle\par}
\vspace{.3cm}{\usebeamerfont{subtitle}\color{upgris}\insertsubtitle\par}
\vspace{.8cm}{\large\insertauthor\par}{\normalsize\insertdate\par}
\end{column}
\begin{column}{.22\textwidth}\vspace{.55cm}\includegraphics[width=\linewidth]{UPlogo.pdf}\end{column}
\end{columns}
\end{frame}

% 2
\begin{frame}{Ruta de la clase}
\begin{enumerate}
\item Problema de Cauchy y equivalencia con una ecuación integral.
\item Métodos de un paso: Euler, trapecio y Heun.
\item Euler mejorado (punto medio) y Runge--Kutta clásico.
\item Comparación del error y aplicación a una trayectoria de capital.
\end{enumerate}
\medskip
Al terminar: deducir las fórmulas, calcular un paso y comprobar qué ocurre al reducir $h$.
\fuente{Secuencia del capítulo 8 de las notas del curso.}
\end{frame}

% 3
\begin{frame}{El problema de Cauchy}
Buscamos una función $\mathbf y(t)\in\R^d$ tal que
\[\boxed{\mathbf y'(t)=\mathbf f(t,\mathbf y(t)),\qquad
\mathbf y(t_0)=\mathbf y_0.}\]
En las notas, $\mathbf f:\Omega\subset\R\times\R^d\to\R^d$ es continua,
con $\Omega$ abierto y convexo.
\begin{itemize}
\item La ecuación fija la velocidad de cambio; el dato inicial selecciona la trayectoria.
\item Continuidad y Lipschitz local en $\mathbf y$ aseguran existencia y unicidad local.
\item Trabajaremos en $[a,b]$, suponiendo que la solución existe en todo el intervalo.
\end{itemize}
\end{frame}

% 4
\begin{frame}{Equivalencia con la ecuación integral}
\begin{block}{Lema del capítulo 8}
Para $\mathbf f$ continua, una función continua $\mathbf y$ satisface
\[\mathbf y(t)=\mathbf y_0+\int_{t_0}^t\mathbf f(s,\mathbf y(s))\,ds\]
si y solo si es solución del problema de Cauchy.
\end{block}
En un paso:
\[\boxed{\mathbf y(t_{k+1})=\mathbf y(t_k)+
\int_{t_k}^{t_{k+1}}\mathbf f(s,\mathbf y(s))\,ds.}\]
La cuadratura del capítulo anterior permite construir métodos para la EDO.
\end{frame}

% 5
\begin{frame}{Demostración del lema}
\textbf{De la EDO a la integral.} Integramos y aplicamos el teorema fundamental:
\[\mathbf y(t)-\mathbf y(t_0)=\int_{t_0}^t\mathbf y'(s)\,ds
=\int_{t_0}^t\mathbf f(s,\mathbf y(s))\,ds.\]
Usamos $\mathbf y(t_0)=\mathbf y_0$.
\medskip

\textbf{De la integral a la EDO.} El integrando es continuo, luego
\[\frac{d}{dt}\left[\mathbf y_0+\int_{t_0}^t\mathbf f(s,\mathbf y(s))\,ds\right]
=\mathbf f(t,\mathbf y(t)).\]
Al sustituir $t=t_0$ recuperamos la condición inicial.
\end{frame}

% 6
\begin{frame}{Malla temporal y métodos de un paso}
Elegimos $n\in\mathbb N$, $h=(b-a)/n$ y
\[t_k=a+kh,\qquad k=0,\ldots,n.\]
En la notación de las notas, $\mathbf y_k\approx\mathbf y(t_k)$ y
\[\boxed{\mathbf y_{k+1}=\mathbf y_k+
h\boldsymbol\phi(t_k,\mathbf y_k;h),\qquad \mathbf y_0=\mathbf y(a).}\]
\begin{itemize}
\item $\boldsymbol\phi$: pendiente efectiva del paso.
\item Un paso usa el estado $\mathbf y_k$; puede evaluar varias pendientes intermedias.
\item En las deducciones siguientes usamos el caso escalar; las fórmulas se aplican a vectores.
\end{itemize}
\end{frame}

% 7
\begin{frame}{Euler: rectángulo izquierdo}
Partimos de la identidad exacta
\[y(t_{k+1})=y(t_k)+\int_{t_k}^{t_k+h}f(s,y(s))\,ds.\]
Aproximamos el integrando por su valor en el extremo izquierdo:
\[\int_{t_k}^{t_k+h}f(s,y(s))\,ds\approx h f(t_k,y(t_k)).\]
Sustituimos el valor exacto por el aproximado:
\[\boxed{y_{k+1}=y_k+h f(t_k,y_k).}\]
Es explícito y necesita una evaluación de $f$ por paso.
\end{frame}

% 8
\begin{frame}{Euler: Taylor e interpretación geométrica}
\begin{columns}[c,onlytextwidth]
\begin{column}{.54\textwidth}
Si $y\in C^2$,
\[y(t_k+h)=y(t_k)+hy'(t_k)
+\frac{h^2}{2}y''(\xi_k).\]
Como $y'=f(t,y)$, Euler conserva los dos primeros términos.
\medskip

Avanza sobre la recta tangente calculada en el estado actual.
\end{column}
\begin{column}{.42\textwidth}
\begin{tikzpicture}[scale=1.05]
\draw[->] (0,0)--(4.4,0) node[right] {$t$};
\draw[->] (0,0)--(0,3.6) node[above] {$y$};
\draw[upazul,thick,domain=.35:3.65,samples=60] plot (\x,{.65+.13*(\x)^2+.2*\x});
\draw[updorado,very thick] (1,.98)--(3,1.90);
\draw[dashed] (1,0) node[below] {$t_k$}--(1,.98);
\draw[dashed] (3,0) node[below] {$t_k+h$}--(3,2.42);
\fill[upazul] (1,.98) circle (2pt);
\fill[updorado] (3,1.9) circle (2pt) node[right] {$y_{k+1}$};
\fill[upazul] (3,2.42) circle (2pt);
\node[upazul] at (2.3,3.2) {solución exacta};
\end{tikzpicture}
\end{column}
\end{columns}
\end{frame}

% 9
\begin{frame}{Ejemplo para toda la clase}
\[y'=y-t^2+1,\qquad y(0)=\tfrac12,\qquad 0\leq t\leq2.\]
La solución exacta, para verificar los cálculos, es
\[y(t)=(t+1)^2-\tfrac12 e^t.\]
Con $h=0.2$, Euler da
\[f(0,0.5)=1.5,\qquad y_1=0.5+0.2(1.5)=\boxed{0.8}.\]
La referencia es $y(0.2)=0.8292986209$.
\medskip

\textbf{En pizarra:} calcule $y_2$. Recuerde evaluar $f$ en $(t_1,y_1)$.
\end{frame}

% 10
\begin{frame}[fragile]{Algoritmo de Euler}
\begin{lstlisting}
import numpy as np

def euler(f, a, b, y0, n):
    h = (b-a)/n             # n entero positivo
    t = np.linspace(a, b, n+1)
    y = np.empty(n+1)
    y[0] = y0
    for k in range(n):
        y[k+1] = y[k] + h*f(t[k], y[k])
    return t, y
\end{lstlisting}
Fijamos $n$ para que la malla llegue exactamente a $b$.
Se actualiza con $y_k$, antes de avanzar al siguiente nodo.
\end{frame}

% 11
\begin{frame}{Error local y error global de Euler}
El \concepto{defecto local} parte del valor exacto en un paso:
\[d_{k+1}=y(t_{k+1})-y(t_k)-hf(t_k,y(t_k))
=\frac{h^2}{2}y''(\xi_k).\]
Si $|y''|\leq M$ y $f$ es Lipschitz en $y$ con constante $L$,
\[|e_{k+1}|\leq(1+hL)|e_k|+\frac M2h^2,
\qquad e_k=y(t_k)-y_k.\]
Con dato inicial exacto y $L>0$,
\[\max_k|e_k|\leq\frac{Mh}{2L}\bigl(e^{L(b-a)}-1\bigr)=O(h).\]
\small Para $L=0$, la cota es $M(b-a)h/2$. El defecto aquí no se divide por $h$.
\end{frame}

% 12
\begin{frame}{Regla del trapecio: fórmula implícita}
La cuadratura trapezoidal sugiere
\[\int_{t_k}^{t_{k+1}}f(s,y(s))\,ds\approx
\frac h2[f(t_k,y(t_k))+f(t_{k+1},y(t_{k+1}))].\]
El método resultante es
\[\boxed{y_{k+1}=y_k+\frac h2[f(t_k,y_k)+f(t_{k+1},y_{k+1})].}\]
\begin{itemize}
\item $y_{k+1}$ aparece también dentro de $f$: en general hay que resolver una ecuación.
\item Si $f$ es lineal en $y$, se puede despejar el nuevo valor.
\item El capítulo 8 realiza una predicción adicional para obtener un método explícito.
\end{itemize}
\end{frame}

% 13
\begin{frame}{Heun: predictor de Euler y corrección}
Primero predecimos el valor final:
\[y_{k+1}^{\mathrm{pred}}=y_k+hf(t_k,y_k).\]
Evaluamos la pendiente en ese punto y promediamos:
\[\boxed{y_{k+1}=y_k+\frac h2\left[f(t_k,y_k)+
f(t_k+h,y_k+hf(t_k,y_k))\right].}\]
\concepto{Heun o trapecio explícito}: dos evaluaciones de $f$ por paso.
\medskip

Una sola corrección con el predictor no resuelve exactamente la ecuación del trapecio implícito.
\end{frame}

% 14
\begin{frame}{Ejemplo: un paso de Heun}
Para $f(t,y)=y-t^2+1$, $t_0=0$, $y_0=0.5$ y $h=0.2$:
\begin{align*}
k_1&=f(0,0.5)=1.5,\\
y_1^{\mathrm{pred}}&=0.5+0.2(1.5)=0.8,\\
k_2&=f(0.2,0.8)=0.8-0.04+1=1.76,\\
y_1&=0.5+\frac{0.2}{2}(1.5+1.76)=\boxed{0.826}.
\end{align*}
La referencia exacta es $0.8292986209$.
\medskip

\textbf{Comparación:} el trapecio implícito da $y_1=0.8288888889$.
\end{frame}

% 15
\begin{frame}{Euler mejorado: punto medio explícito}
Las notas llaman \concepto{Euler mejorado} al siguiente RK de dos etapas:
\[k_1=f(t_k,y_k),\qquad
k_2=f\!\left(t_k+\frac h2,y_k+\frac h2k_1\right),\]
\[\boxed{y_{k+1}=y_k+h k_2.}\]
Para el mismo primer paso,
\[k_2=f(0.1,0.65)=1.64,\qquad y_1=\boxed{0.828}.\]
La pendiente del punto medio se evalúa en un estado aproximado.
\medskip

Usaremos el nombre \textbf{punto medio explícito} para identificar la fórmula.
\end{frame}

% 16
\begin{frame}{Dos métodos de orden dos}
\begin{center}
\begin{tabular}{lcc}\toprule
&Punto medio&Heun\\\midrule
Tiempo de la segunda etapa&$t_k+h/2$&$t_k+h$\\[2mm]
Pesos de $k_1,k_2$&$0,1$&$1/2,1/2$\\[2mm]
Primer paso del ejemplo&$0.828$&$0.826$\\[2mm]
Evaluaciones por paso&2&2\\\bottomrule
\end{tabular}
\end{center}
Ambos tienen defecto local $O(h^3)$ y error global $O(h^2)$, bajo suficiente suavidad y Lipschitz.
\medskip

El mismo orden no implica el mismo error ni la misma trayectoria.
\end{frame}

% 17
\begin{frame}{Runge--Kutta explícito: forma general}
Con la notación $a_i,b_{i\ell},c_i$ del capítulo 8:
\begin{align*}
k_1&=f(t_k,y_k),\\
k_i&=f\!\left(t_k+a_i h,\ y_k+h\sum_{\ell=1}^{i-1}b_{i\ell}k_\ell\right),
\quad i=2,\ldots,m,\\
\phi(t_k,y_k;h)&=\sum_{i=1}^m c_i k_i,\\
y_{k+1}&=y_k+h\sum_{i=1}^m c_i k_i.
\end{align*}
Cada etapa usa solamente pendientes ya calculadas.
Los $k_i$ son pendientes: el factor $h$ aparece en los incrementos.
\end{frame}

% 18
\begin{frame}{Coeficientes y consistencia}
Se suele imponer
\[a_1=0,\qquad a_i=\sum_{\ell<i}b_{i\ell}.\]
Así, las etapas usan desplazamientos temporales y de estado compatibles.
\medskip

Para integrar exactamente $y'=1$, todas las pendientes son $k_i=1$ y
\[y_{k+1}=y_k+h\sum_i c_i.\]
Por tanto, necesitamos \concepto{$\sum_i c_i=1$}.
Para orden dos se exige además $\sum_i c_i a_i=1/2$.
\medskip

El número de etapas por sí solo no determina el orden del método.
\end{frame}

% 19
\begin{frame}{Tablas de coeficientes de las notas}
\begin{columns}[T,onlytextwidth]
\begin{column}{.46\textwidth}
\centering\textbf{Punto medio}
\[\begin{array}{c|cc}0&0&0\\1/2&1/2&0\\\hline &0&1\end{array}\]
\textbf{Heun}
\[\begin{array}{c|cc}0&0&0\\1&1&0\\\hline &1/2&1/2\end{array}\]
\end{column}
\begin{column}{.48\textwidth}
\centering\textbf{RK4 clásico}
\[\begin{array}{c|cccc}
0&0&0&0&0\\
1/2&1/2&0&0&0\\
1/2&0&1/2&0&0\\
1&0&0&1&0\\\hline
&1/6&1/3&1/3&1/6
\end{array}\]
\raggedright\small Columna izquierda: $a_i$. Interior: $b_{i\ell}$. Fila inferior: $c_i$.
\end{column}
\end{columns}
\medskip
\textbf{En pizarra:} compruebe $\sum c_i=1$ y $\sum c_i a_i=1/2$.
\end{frame}

% 20
\begin{frame}{Runge--Kutta clásico de cuarto orden}
\begin{align*}
k_1&=f(t_k,y_k),\\
k_2&=f(t_k+h/2,y_k+hk_1/2),\\
k_3&=f(t_k+h/2,y_k+hk_2/2),\\
k_4&=f(t_k+h,y_k+hk_3).
\end{align*}
\[\boxed{y_{k+1}=y_k+\frac h6(k_1+2k_2+2k_3+k_4).}\]
Cuatro evaluaciones por paso. Bajo suficiente suavidad y Lipschitz,
el defecto local es $O(h^5)$ y el error global es $O(h^4)$.
\end{frame}

% 21
\begin{frame}{Ejemplo: las cuatro etapas de RK4}
Con $h=0.2$ y $(t_0,y_0)=(0,0.5)$:
\begin{align*}
k_1&=1.5,\\
k_2&=f(0.1,0.65)=1.64,\\
k_3&=f(0.1,0.664)=1.654,\\
k_4&=f(0.2,0.8308)=1.7908.
\end{align*}
\[y_1=0.5+\frac{0.2}{6}(1.5+2(1.64)+2(1.654)+1.7908)
=\boxed{0.8292933333}.\]
Error absoluto del primer paso: $5.29\times10^{-6}$.
\end{frame}

% 22
\begin{frame}{Orden y costo por trayectoria}
\begin{center}
\begin{tabular}{lccc}\toprule
Método&Defecto local&Error global&Evaluaciones\\\midrule
Euler&$O(h^2)$&$O(h)$&$n$\\[1mm]
Punto medio&$O(h^3)$&$O(h^2)$&$2n$\\[1mm]
Heun&$O(h^3)$&$O(h^2)$&$2n$\\[1mm]
RK4&$O(h^5)$&$O(h^4)$&$4n$\\\bottomrule
\end{tabular}
\end{center}
Órdenes para una solución suficientemente suave, en un intervalo fijo.
\medskip

La comparación debe considerar el paso, la precisión requerida y el costo de evaluar $f$.
Una etapa adicional no exige derivar $f$ analíticamente.
\end{frame}

% 23
\begin{frame}{Comparación numérica con solución exacta}
\[E_h=\max_{0\leq k\leq n}|y(t_k)-y_k|,\qquad h=2/n.\]
\begin{center}\small\input{\datos/tabla_errores.tex}\end{center}
Se usa el mismo PVI $y'=y-t^2+1$, $y(0)=0.5$, hasta $t=2$.
\medskip

Son errores máximos sobre la malla; no constituyen una cota entre nodos.
\end{frame}

% 24
\begin{frame}{Refinamiento y orden observado}
Si $E_h\approx Ch^p$, entonces
\[\frac{E_h}{E_{h/2}}\approx2^p,\qquad
p_{\mathrm{obs}}=\log_2\!\left(\frac{E_h}{E_{h/2}}\right).\]
\begin{center}\input{\datos/tabla_ordenes.tex}\end{center}
\small Se comparan $n=40$ y $n=80$.
Con mallas demasiado gruesas o errores próximos al redondeo, el cociente puede no mostrar el orden teórico.
\end{frame}

% 25
\begin{frame}{Trayectorias en una misma malla}
\begin{center}
\begin{tikzpicture}
\begin{axis}[width=.88\textwidth,height=.67\textheight,
xlabel={$t$},ylabel={$y$},xmin=0,xmax=2,grid=major,
legend style={font=\footnotesize,at={(.03,.97)},anchor=north west},
tick label style={font=\small},label style={font=\small}]
\addplot[black,thick,domain=0:2,samples=100] {(x+1)^2-.5*exp(x)};
\addlegendentry{Exacta}
\addplot[uprojo,mark=o] table[x=t,y=Euler] {\datos/trayectorias.dat};
\addlegendentry{Euler}
\addplot[upverde,mark=triangle*] table[x=t,y=Heun] {\datos/trayectorias.dat};
\addlegendentry{Heun}
\addplot[upazul,only marks,mark=+,mark size=3] table[x=t,y=RK4] {\datos/trayectorias.dat};
\addlegendentry{RK4}
\end{axis}
\end{tikzpicture}
\end{center}
\small $n=10$, $h=0.2$. Las tablas permiten comparar diferencias poco visibles en la figura.
\end{frame}

% 26
\begin{frame}{Actividad: 80 evaluaciones de la función}
Suponga que dispone de \textbf{80 evaluaciones de $f$} para llegar a $t=2$.
\begin{enumerate}
\item Determine $n$ y $h$ para Euler, punto medio, Heun y RK4.
\item Use las tablas o el cuaderno para comparar los errores.
\item Explique por qué punto medio y Heun pueden producir errores distintos.
\end{enumerate}
\medskip
Después, sin usar la solución exacta, proponga una comparación entre dos mallas para estimar el error final.
\medskip
\small Trabaje primero en parejas; contraste los resultados con el cuaderno.
\end{frame}

% 27
\begin{frame}{Aplicación: acumulación de capital}
Retomamos el modelo hipotético de la clase de integración:
\[K'(t)=I(t)-\delta K(t),\qquad K(0)=10,\]
\[I(t)=2e^{0.02t},\qquad\delta=0.05,\qquad0\leq t\leq10.\]
$K$ se mide en unidades de capital, $t$ en años e $I$ en unidades de capital por año.
La referencia exacta es
\[K(t)=10e^{-0.05t}+\frac{2}{0.07}\left(e^{0.02t}-e^{-0.05t}\right).\]
No cambia la estructura del método: evaluamos $f(t,K)=I(t)-\delta K$ en cada etapa.
\end{frame}

% 28
\begin{frame}{Capital: Heun y comparación de resultados}
El predictor y la corrección son
\[K_{k+1}^{\mathrm{pred}}=K_k+h[I(t_k)-\delta K_k],\]
\[K_{k+1}=K_k+\frac h2\bigl[I(t_k)-\delta K_k+
I(t_{k+1})-\delta K_{k+1}^{\mathrm{pred}}\bigr].\]
Con $n=20$ y $h=0.5$, $K(10)=23.6330808385$:
\begin{center}\small
\begin{tabular}{lrr}\toprule
Método&Aproximación de $K(10)$&Error absoluto\\\midrule
Euler&23.6791486574&$4.607\times10^{-2}$\\
Punto medio&23.6326201571&$4.607\times10^{-4}$\\
Heun&23.6328425105&$2.383\times10^{-4}$\\
RK4&23.6330808293&$9.137\times10^{-9}$\\\bottomrule
\end{tabular}\end{center}
\end{frame}

% 29
\begin{frame}{Sistemas: las mismas fórmulas, con vectores}
Una ecuación de segundo orden también puede escribirse como sistema:
\[u''=-u,\quad u(0)=1,\quad u'(0)=0,\qquad
\mathbf y=\begin{pmatrix}u\\v\end{pmatrix},\quad
\mathbf f(t,\mathbf y)=\begin{pmatrix}v\\-u\end{pmatrix}.\]
Con Euler y $h=0.1$,
\[\mathbf y_1=\begin{pmatrix}1\\0\end{pmatrix}
+0.1\begin{pmatrix}0\\-1\end{pmatrix}
=\begin{pmatrix}1\\-0.1\end{pmatrix}.\]
Todas las componentes se actualizan desde el \textbf{mismo estado anterior}.
En RK4, cada $\mathbf k_i$ es un vector de pendientes.
\end{frame}

% 30
\begin{frame}{Complemento: paso y estabilidad}
Para $y'=\lambda y$, $\lambda<0$, la solución exacta decae y
\[y_{k+1}=R(z)y_k,\qquad z=h\lambda.\]
\[R_{\mathrm E}(z)=1+z,\qquad
R_{\mathrm H}(z)=1+z+z^2/2.\]
Si $|R(z)|>1$, las perturbaciones crecen de un paso al siguiente.
\medskip

Para $y'=-4y$, con $h=0.75$, Euler tiene $R=-2$ y Heun $R=2.5$.
Ambos amplifican, aunque Heun sea de orden dos.
\medskip

\small En el eje real negativo, ambos satisfacen $|R(z)|<1$ para $-2<z<0$.
Este diagnóstico amplía el contenido básico del capítulo.
\end{frame}

% 31
\begin{frame}{Preguntas de cierre}
\begin{enumerate}
\item Obtenga Euler y Heun a partir de la identidad integral de un paso.
\item Distinga el predictor $y_{k+1}^{\mathrm{pred}}$ del valor corregido.
\item Escriba las cuatro etapas de RK4 e identifique dónde interviene $h$.
\item Si al dividir $h$ por dos el error baja aproximadamente por un factor 16,
\textquestiondown qué orden sugiere el experimento?
\end{enumerate}
\medskip
Para verificar una trayectoria: condición inicial, fórmula de paso, refinamiento y error respecto de una referencia disponible.
\end{frame}

% 32
\begin{frame}{Lectura y trabajo posterior}
\begin{itemize}
\item \textbf{Notas del curso, capítulo 8}: problema de Cauchy, métodos de un paso,
Euler, trapecio y Runge--Kutta explícitos.
\item \textbf{Capítulo 7}: reglas de cuadratura utilizadas en la construcción.
\item \textbf{Cuaderno de esta clase}: pasos manuales, refinamiento, capital y sistemas.
\end{itemize}
\medskip
Trabajo: repita la comparación para $y'=y$, $y(0)=1$, hasta $t=2$.
Explique por qué punto medio y Heun coinciden en este problema, aunque sus fórmulas sean diferentes.
\end{frame}
\end{document}
